\paragraph{Problem 4 answer}~\GradescopeComment{This should be on page 8 of your PDF.}

Fix the input size $N$.

\begin{enumerate}[label=(\roman*)]
\item[(i)] For $n\le N$, define $T(n)$ to be run-time on sub-arrays of size $n$.

  Then, for $n\le \sqrt N$, we have $T(n) = \BLUE{
    \REPLACEME 
  }$ 

  For $n>\sqrt N$, we have $T(n) = \BLUE{
    \REPLACEME 
  }$
  
\item[(ii)] Within each level $i$:

  The number of nodes is $\BLUE{
    \REPLACEME 
  }$
  
  The problem size is $\BLUE{
    \REPLACEME 
  }$
  
  If $i$ is not the bottom level, the work done per subproblem is $\BLUE{
    \REPLACEME 
  }$

  If $i$ is the bottom level, the work done per subproblem is $\BLUE{
    \REPLACEME 
  }$
  %
  % ADD AN EXPLANATION IF YOU WANT
  %
  
\item[(iii)] The number of levels is $\BLUE{
    %
    \REPLACEME
    %
  }$, because \BLUE{
    %
    \REPLACEME
    % 
  .}
  
\item[(iv)] Summing over the levels, the total work is
  \[\BLUE{ T(N) ~=~  
      %
      \REPLACEME
      % 
    .}\]

\item[(v)] The big-$\Theta$ value of this sum is
  \[\BLUE{
      %
      \REPLACEME
      % 
    .}\]

\end{enumerate}


\vfill

\paragraph{Problem \arabic{problem} collaboration and citations}~\GradescopeComment{This should be on page 8 (bottom) of your PDF.}

I collaborated on this problem with the following people:
\begin{blue}
  \begin{itemize}
  \item \REPLACEME                    % write "none" if none
  \end{itemize}
\end{blue}

My answers use ideas from the following sources (other than the lecture notes and textbooks): 
\begin{blue}
  \begin{itemize}
  \item \REPLACEME                    % write "none" if none
  \end{itemize}
\end{blue}
 
%%% Local Variables:
%%% mode: latex
%%% TeX-master: "homework_2"
%%% End:
